% Hafta 6 — Cihaz bağlama: kopyalama işe yaramasın
% Derleme: py -3.12 tools/latex/derle.py docs/week-6/assets/h06-12-cihaz-baglama.tex
\documentclass[tikz,border=8pt]{standalone}
\usepackage{cen429-cizim}
\begin{document}
\begin{tikzpicture}
  \node[baslik] (b) at (0,0) {Cihaz bağlama: kopyalama işe yaramasın};
  \node[altbaslik] at (b.south west) {çalınan dosya başka cihazda anlamsızdır};
  \node[kutu=36mm, anchor=north west] (n1) at (0,-2.4)
    {\kb{Cihaz parmak izi}\\donanım/OS özellikleri};
  \node[koyukutu=34mm, right=9mm of n1] (n2)
    {\kb{HKDF ile türet}\\cihaza özgü anahtar};
  \node[kutu=34mm, right=9mm of n2] (n3)
    {\kb{Sırrı bu anahtarla sar}\\AEAD};
  \node[iyikutu=36mm, right=9mm of n3] (n4)
    {\kb{Başka cihazda}\\parmak izi farklı $\rightarrow$ AÇILMAZ};
  \draw[ok] (n1.east) -- (n2.west);
  \draw[ok] (n2.east) -- (n3.west);
  \draw[iyiok] (n3.east) -- (n4.west);
  \node[fit=(n1)(n2)(n3)(n4), inner sep=0pt] (grp) {};
  \node[dip, text width=155mm, anchor=north west] at ($(grp.south west)+(0,-8mm)$)
    {Code lifting ve kopyalama saldırılarının ölçeklenmesini bu engeller (11. hafta).};
\end{tikzpicture}
\end{document}
